\documentclass[11pt]{book}
\usepackage{fontspec}
\setmainfont{Latin Modern Roman}
\setsansfont{Latin Modern Sans}
\usepackage{amsmath,amssymb,mathtools}
\usepackage{unicode-math}
\setmathfont{Latin Modern Math}
\usepackage{xurl}
\usepackage[hidelinks]{hyperref}
\newcommand{\sref}[2]{\hyperlink{source:#1:#2}{[S#2]}}
\newcommand{\cataloguescope}[1]{\paragraph{Mathematical question.}#1}
\begin{document}
\section{Dirac's Conjecture}\label{statement}

\subsection{Definitions and mathematical statement}

\paragraph{Shared notation.}$\mathbb N=\{1,2,\ldots\}$, and $\mathbb Z,\mathbb Q,\mathbb R,\mathbb C$ denote the integers, rationals, reals and complex numbers. Any other conventions are specified locally.


For a finite non-collinear set $P\subset\mathbb R^2$ with $n=|P|\geq3$, let $\mathcal L(P)$ be the set of distinct affine lines containing at least two points of $P$, and let $d_P(p)=|\{\ell\in\mathcal L(P):p\in\ell\}|$.
\paragraph{Statement recorded in the supplied source.}
Is there an absolute real constant $c\geq0$, independent of $n$ and $P$, such that
\[\forall P\subset\mathbb R^2\text{ finite and non-collinear},\qquad\max_{p\in P}d_P(p)\geq \frac{|P|}{2}-c?\] \sref{statement}{2}
\subsection{Short English statement}
Must some point of every non-collinear planar configuration lie on at least half as many determined lines as there are points, up to one fixed additive constant?
\subsection{Sources}
\begin{description}
\item[\hypertarget{source:statement:1}{S1}] ULAM.AI and the UnsolvedMath Contributors, \emph{UnsolvedMath: A Curated Collection of Open Mathematics Problems}, catalogue entry. CC BY 4.0. \url{https://huggingface.co/datasets/ulamai/UnsolvedMath}.
\item[\hypertarget{source:statement:2}{S2}] Dirac, Gabriel, Dirac's Conjecture, Open Problem Garden, node 59888. Complete problem, discussion and bibliography (including mathematical image alt text). \url{https://www.openproblemgarden.org/op/diracs_conjecture}.
\end{description}
\label{end:statement}
\end{document}
